%=======================================================================
%  Ejercicio 6 — Vuelos de Lévy en el modelo de hormigas forrajeras (Tipo 2)
%=======================================================================
\setcounter{ejercicio}{5}   % el enunciado lo numera 6
\begin{ejercicio}[Tipo 2]\label{ej:levy}
\textbf{Aplicación de vuelos de Lévy en el modelo de hormigas forrajeras.}
En este ejercicio, explorarás una estrategia de movimiento alternativa para las
hormigas en el modelo de «Ants» de la biblioteca de NetLogo. Normalmente, las
hormigas realizan caminatas aleatorias para buscar comida. Modificarás este
comportamiento implementando los \emph{vuelos de Lévy}, un patrón de movimiento
que se utiliza frecuentemente en biología para describir estrategias de búsqueda
eficientes.

\paragraph{Parte 1: Entendiendo el modelo de hormigas}
\begin{enumerate}
  \item \textbf{Explora el modelo existente:} Ejecuta el modelo de «Ants» de la
    biblioteca de NetLogo para observar el comportamiento predeterminado de las
    hormigas mientras buscan comida. Presta atención a cómo se mueven y
    responden a los rastros de feromonas. Siéntete libre de modificar el
    contexto, mira cómo se comporta en los límites.
  \item \textbf{Comportamiento de caminata aleatoria:} Analiza el movimiento de
    las hormigas en ausencia de comida. Demuestra a través de observaciones que,
    cuando no hay comida ni feromonas presentes, las hormigas realizan una
    caminata aleatoria.
    \begin{itemize}
      \item Mide la distribución de las longitudes de paso y los ángulos de giro
        a lo largo del tiempo.
      \item Confirma que, sin estímulos externos (como comida o feromonas), el
        movimiento de las hormigas coincide con las características de una
        caminata aleatoria.
    \end{itemize}

    Revisando el código, vi que la sección que se encarga de la aleatoriedad del movimiento de las hormigas es la siguiente:
    \begin{lstlisting}
      to wiggle  ;; turtle procedure
        rt random 40
        lt random 40
        if not can-move? 1 [ rt 180 ]
      end
      \end{lstlisting}

      En resumen, lo que hace la función es buscar dos número aleatorios entre $0$ y $40$, y girar hacia la diferencia.
      Esto de cierta forma podría representar una caminata aleatoria pues la dirección a la que van las hormigas si es 
      "aleatoria", sin embargo, es una caminata aleatoria sesgada, pues la dirección a la que van las hormigas sigue una tendencia
      a ir derecho. Podemos hacer un histograma de los angulos de giro de las hormigas y ver que el ángulo más frecuente es $0$, y el
      histograma sigue una forma triangular. En una caminata aleatoria normal, no debería de haber mucha difrencia entre todos 
      los ángulos, y deberían ser igualmente frecuentes.
        
\end{enumerate}

\paragraph{Parte 2: Introducción a los vuelos de Lévy}
Los vuelos de Lévy se caracterizan por pasos cortos intercalados con pasos
largos; los cortos son más frecuentes que los largos, lo que permite una
exploración más eficiente en entornos escasos.
\begin{enumerate}
  \item \textbf{Descripción matemática:} Un vuelo de Lévy se modela con
    longitudes de paso que siguen una distribución de ley de potencias:
    \[
      P(l) \propto l^{-\mu},
    \]
    donde $l$ es la longitud del paso y $\mu$ es el exponente que controla la
    distribución. Típicamente, $1<\mu\le 3$ para sistemas biológicos. Realiza
    las gráficas para $P(l)$ con tres valores de $\mu$. ¿Qué pasa conforme el
    valor crece? ¿Cómo lo interpretas si haces una caminata por CU?
  \item \textbf{Beneficios de los vuelos de Lévy:} Discute la eficiencia de los
    vuelos de Lévy para buscar recursos (como comida) escasos, en comparación
    con las caminatas aleatorias. ¿A qué animales podría beneficiar? ¿A los que
    vuelan, nadan o corren?
\end{enumerate}

\paragraph{Parte 3: Modificación del movimiento de las hormigas}
\begin{enumerate}
  \item \textbf{Implementación de vuelos de Lévy:} Modifica el modelo de
    NetLogo para que las hormigas sigan un patrón de vuelo de Lévy en lugar de
    una caminata aleatoria simple. Ajusta el código para generar longitudes de
    paso de acuerdo con una distribución de ley de potencias.
  \item \textbf{Simulación y comparación:} Ejecuta el modelo modificado y
    observa las diferencias en la eficiencia de búsqueda de comida entre las
    hormigas que usan vuelos de Lévy y aquellas que usan la caminata aleatoria
    predeterminada.
\end{enumerate}

\paragraph{Parte 4: Análisis y discusión}
\begin{enumerate}
  \item \textbf{Análisis de eficiencia:} Mide y compara el tiempo que tardan
    las hormigas en encontrar comida usando ambos patrones de movimiento (vuelos
    de Lévy vs.\ caminata aleatoria). Recuerda que estás realizando simulaciones
    aleatorias, una simulación no basta, debes poder medir la media y su
    desviación estándar.
  \item ¿En qué condiciones ambientales los vuelos de Lévy superan a las
    caminatas aleatorias?
\end{enumerate}
Añade tu código en un archivo separado.
\end{ejercicio}

\begin{solucion}

\paragraph{Parte 1}
\begin{enumerate}
  \item \pendiente{}
  \item \pendiente{}
\end{enumerate}

\paragraph{Parte 2}
\begin{enumerate}
  \item \pendiente{} \parencite{lakhloufi2025power}
  \item \pendiente{}
\end{enumerate}

\paragraph{Parte 3}
\begin{enumerate}
  \item \pendiente{}
  \item \pendiente{}
\end{enumerate}

\paragraph{Parte 4}
\begin{enumerate}
  \item \pendiente{}
  \item \pendiente{}
\end{enumerate}

\paragraph{Código}
\pendiente{nombre del archivo con el modelo modificado}
\end{solucion}
